%% ==============================
\Appendix
\label{ch:Appendix}
%% ==============================

\section{Reproduzierbarkeit des Messaufbaus}
\label{sec:app-reproduzierbarkeit}

Dieser Anhang enthält diejenigen Angaben zum Messaufbau, die für eine Nachbildung der Messungen erforderlich, für den Argumentationsgang der Arbeit jedoch nicht tragend sind.

\begin{table}[htbp]
\renewcommand{\arraystretch}{1.3}
\footnotesize
  \centering
  \caption{Ausführungsumgebung des Messaufbaus}
  \label{tab:impl-umgebung}
  \begin{tabular}{p{0.34\textwidth} p{0.56\textwidth}}
    \toprule
    \textbf{Merkmal} & \textbf{Wert} \\
    \midrule
    Prozessor & AMD Ryzen 7 PRO 8840HS, 3,3\,GHz Grundtakt, 8 Kerne / 16 Threads \\
    Prozessorarchitektur & x86-64 \\
    Arbeitsspeicher des Wirtssystems & 32\,GB (physisch), 27,7\,GB nutzbar \\
    Der Containerlaufzeit zugewiesen (WSL2) & 14 Prozessorkerne, 21,5\,GiB Arbeitsspeicher \\
    Betriebssystem & Windows 11 Enterprise, Version 10.0.26100, Build 26100 \\
    Containerlaufzeit & Docker 29.7.2 (API 1.55), Client 29.6.2-rd \\
    Linux-Subsystem & WSL 2.7.14.0, Kernel 6.18.33.2-2 \\
    \bottomrule
  \end{tabular}
\end{table}
\normalsize


\begin{table}[htbp]
\renewcommand{\arraystretch}{1.3}
\footnotesize
  \centering
  \caption{Startparameter des Agenten je Variante. Der Startwert der Trustee-DID ist gesetzt, wird hier jedoch nicht abgedruckt}
  \label{tab:impl-startparameter}
  \begin{tabular}{p{0.26\textwidth} p{0.31\textwidth} p{0.31\textwidth}}
    \toprule
    \textbf{Parameter} & \textbf{HCS-Variante} & \textbf{Indy-Variante} \\
    \midrule
    \multicolumn{3}{l}{\textit{übereinstimmend}} \\
    Eingehender Transport & \texttt{http 0.0.0.0 3000} & identisch \\
    Ausgehender Transport & \texttt{http} & identisch \\
    Endpunkt & \texttt{http://localhost:3000} & identisch \\
    Verwaltungsschnittstelle & \texttt{0.0.0.0 3001}, ungesicherter Modus & identisch \\
    Wallet-Typ & \texttt{askar-anoncreds} & identisch \\
    Selbsteinrichtung & aktiviert & identisch \\
    Protokollierungsstufe & \texttt{info} & identisch \\
    \midrule
    \multicolumn{3}{l}{\textit{durch die Ledger-Anbindung erzwungen}} \\
    Registry-Anbindung & \texttt{-{}-no-ledger}, Hedera-Erweiterung geladen & Genesis-Datei des Netzwerks \\
    Bootstrap-Identität & keine; DID wird zur Laufzeit über den Registrierungsendpunkt erzeugt & Startwert gesetzt (Wert nicht abgedruckt), unsicherer Startwert zugelassen \\
    Zugangsdaten & Betreiberdatei mit Konto- und Schlüsselangaben & entfällt \\
    \midrule
    \multicolumn{3}{l}{\textit{übrige Unterschiede}} \\
    Bezeichnung und Wallet-Name & \texttt{bare-vdr-hcs}, \texttt{hcs\_arm} & \texttt{bare-vdr-indy}, \texttt{indy\_arm} \\
    Protokollierung der Ledger-Bibliothek & nicht gesetzt & \texttt{debug} für \texttt{indy-vdr} \\
    \bottomrule
  \end{tabular}
\end{table}
\normalsize



\begin{table}[htbp]
\renewcommand{\arraystretch}{1.3}
\footnotesize
  \centering
  \caption{Parameter der Szenario-Module}
  \label{tab:impl-szenarioparameter}
  \begin{tabular}{p{0.08\textwidth} p{0.80\textwidth}}
    \toprule
    \textbf{Szenario} & \textbf{Parameter} \\
    \midrule
    S1 & Serienlänge 15 auswertbare Läufe je Variante und Messpunkt; erster Lauf als Funktionsnachweis ausgeschlossen. Schreib- und Lesewerte in getrennten Abläufen erhoben \\
    Grundlast & Getakteter Nachrichtenstrom auf einem eigenen Topic, Taktintervall 3,0\,s, über die gesamte Messdauer aufrechterhalten \\
    S2 & Startstufe 5 parallele Nutzer, Schrittweite 5, höchstens 20 Stufen. Je Stufe 60\,s Anlauf (nicht gemessen) und 120\,s Messfenster. Abbruchkriterien: Durchsatzzuwachs unter 5\,\%, Antwortzeit im 95.~Perzentil über dem Dreifachen der Ausgangsantwortzeit, Fehlerrate über 5\,\% \\
    S3 & Zwischenstände bei 10, 30, 100 und 300 Füllobjekten; auf der Indy-Variante zusätzlich bis 2700. Je Zwischenstand drei Lesungen des Zielobjekts und eine Kontrolllesung \\
    S4 & Ausfall nach 120\,s Vorlauf, Ausfallfenster 180\,s, Beobachtung der Erholung über 300\,s bzw. 600\,s. Zielknoten: Knoten~2 auf der Indy-Variante, Knoten~1 auf der HCS-Variante -- dort allerdings ohne Wirkung der Angabe, da die Klientbibliothek unabhängig davon stets denselben Knoten anspricht (Abschnitt~\ref{sec:impl-befunde}, Befund B5). Fünf Wiederholungen je Variante \\
    S5 & Einzelnachweise ohne Serie, ausgeführt über direkte Ledger-Aufrufe unterhalb des Agenten \\
    \bottomrule
  \end{tabular}
\end{table}
\normalsize